% =============================================================================
% cap04.tex — DEFINICAO DO PROBLEMA
% =============================================================================

\chapter{DEFINIÇÃO DO PROBLEMA}\label{cap:definicao-problema}

Este capítulo define a estrutura a ser dimensionada e as condições sob as quais
o dimensionamento é feito. Descreve a edificação e a nomenclatura dos seus
elementos, os materiais adotados e os valores característicos declarados como
dados de entrada, e registra as premissas assumidas, incluindo as reservas
normativas que condicionam a validade dos resultados.

As tabelas deste capítulo contêm apenas dados de entrada, isto é, valores
fixados antes da análise e que o programa de cálculo recebe sem os calcular. As
dimensões que resultam do cálculo, como a altura e a largura da viga e as seções
dos pilares, e todos os esforços e taxas de utilização são apresentados em
tabelas próprias de resultados, a partir do \autoref{cap:predimensionamento}.

\section{DESCRIÇÃO DA EDIFICAÇÃO}\label{sec:descricao-edificacao}

Trata-se de uma cobertura destinada a estacionamento de veículos, aberta nas
quatro faces, com 40,5 m de largura por 46,0 m de comprimento, perfazendo
1.863 m\textsuperscript{2} em planta. A estrutura é composta por nove pórticos
transversais igualmente espaçados a cada 5,75 m.

Cada pórtico é formado por três pilares, sendo dois laterais e um central, e
por duas vigas de duplo-bisel de 20,25 m de vão cada uma, apoiadas sobre eles.
Os pilares têm 7,0 m de altura. Como cada viga tem a cumeeira a meio vão, a
cobertura resulta em duas águas duplas dispostas lado a lado, com calha central
sobre a linha de pilares intermediários.

A opção por dois vãos de 20,25 m em vez de um único vão livre de 40,5 m decorre
do uso previsto: em um estacionamento, uma linha de pilares intermediários é
compatível com o arranjo das vagas e da circulação, e evita a altura de viga
que um vão livre dessa ordem exigiria. A \autoref{fig:modelo-3d} apresenta o
conjunto e a \autoref{fig:portico-elevacao}, a elevação de um pórtico. Ambas
provêm do modelo de elementos finitos com seção transversal real, gerado a
partir da geometria de projeto, em que vigas, pilares e terças são de madeira
lamelada colada; os sistemas de contraventamento e a telha não são
representados.

\figura
{Vista geral do modelo tridimensional da edificação}
{0.98}
{fig/modelo_3d.png}
{Elaborado pelo autor (2026)}
{modelo-3d}
{}
{}

A disposição das duas águas duplas lado a lado, com a calha central sobre a
linha de pilares intermediários, aparece com clareza na elevação: são dois vãos
de 20,25 m sobre três pilares de 7,0 m, com a altura da viga crescendo de
1,00 m no apoio a 2,785 m na cumeeira de cada vão, e as peças menores sobre o
extradorso são as terças.

\figura
{Elevação de um pórtico transversal}
{0.98}
{fig/portico_elevacao.png}
{Elaborado pelo autor (2026)}
{portico-elevacao}
{}
{}

A \autoref{tab:geometria-edificacao} reúne os dados geométricos de entrada,
fixados pelo uso da edificação. A inclinação, as alturas e a largura da viga e
as seções dos pilares não constam dela, porque resultam do pré-dimensionamento e
estão reunidas na \autoref{tab:criterios-predim}.

\tabela{Dados geométricos de entrada da edificação}{\footnotesize
\begin{tabular}{lcrl}
\toprule
    \textbf{Parâmetro} & \textbf{Símbolo} & \textbf{Valor} & \textbf{Unidade} \\
\midrule
    Largura total em planta              & $B_{tot}$ & 40,50 & m \\
    Comprimento total em planta          & $L_{tot}$ & 46,00 & m \\
    Vão de cada pórtico                  & $L$       & 20,25 & m \\
    Número de vãos por pórtico           & --        & 2     & -- \\
    Número de pórticos                   & --        & 9     & -- \\
    Espaçamento entre pórticos           & $L_T$     &  5,75 & m \\
    Altura dos pilares                   & $H$       &  7,00 & m \\
    Seção das terças                     & $b_t \times h_t$ & 0,20 $\times$ 0,58 & m \\
    Espaçamento entre terças             & $s$       &  2,53 & m \\
\bottomrule
\end{tabular}}{Elaborado pelo autor (2026)}{geometria-edificacao}{}{}

\section{NOMENCLATURA DOS ELEMENTOS}\label{sec:nomenclatura}

Os elementos recebem nomes que se repetem nos desenhos e em todo o memorial. A
locação usa duas famílias de eixos: os eixos 1 a 9 marcam os pórticos, a cada
5,75 m ao longo do comprimento, e os alinhamentos A, B e C marcam as linhas de
pilares, a cada 20,25 m ao longo da largura, sendo B a linha central. Os
pilares são numerados por pórtico, de A para C, e as vigas por pórtico, primeiro
o vão A--B e depois o vão B--C. A \autoref{tab:nomenclatura} resume o esquema,
e a \autoref{fig:des-planta-cobertura} mostra a planta de cobertura com os
nomes. Os demais desenhos constam do \autoref{ap:desenhos} e são reproduzidos no
texto onde cada elemento é discutido.

\tabela{Nomenclatura dos elementos}{\footnotesize
\begin{tabular}{p{2.5cm}p{3.5cm}p{7.0cm}}
\toprule
    \textbf{Nome} & \textbf{Elemento} & \textbf{Regra ou posição} \\
\midrule
    1 a 9, A a C & Eixos & Pórticos ao longo do comprimento; alinhamentos de pilares ao longo da largura \\
    V1 a V18 & Vigas & V$(2i-1)$ no vão A--B e V$(2i)$ no vão B--C do pórtico $i$ \\
    P1 a P27 & Pilares & P$(3i-2)$ em A, P$(3i-1)$ em B e P$(3i)$ em C, no pórtico $i$; cantos P1, P3, P25 e P27 \\
    PF1, PF2 & Pilares de frontão & Eixo 1 a 5,75 m de A; eixo 9 a 5,75 m de C \\
    T1 a T17 & Linhas de terças & De A para C, a cada 2,53 m; T1, T9 e T17 sobre os apoios \\
    CV-H1, CV-H2 & Contraventamento da cobertura & Baías 1--2 e 8--9 \\
    CV-V1 a CV-V8 & Painéis verticais & A: 1--2, 8--9; C: 1--2, 8--9; B: 1--2, 2--3, 7--8, 8--9 \\
    CV-T1, CV-T2 & Painéis dos frontões & P1--PF1 no eixo 1; PF2--P27 no eixo 9 \\
    L1, L1c, L1f & Ligações de topo & Viga sobre pilar lateral, sela do pilar central e topo do pilar de frontão \\
    L3, L3c, L3f, L3a & Ligações de base & Corrente lateral, corrente central, pilar de frontão e base ancorada com chaveta \\
\bottomrule
\end{tabular}}{Elaborado pelo autor (2026)}{nomenclatura}{}{}

\figura
{Planta de cobertura e locação}
{1.0}
{fig/des01_planta_cobertura.pdf}
{Elaborado pelo autor (2026)}
{des-planta-cobertura}
{}
{}

\section{MATERIAIS}\label{sec:materiais}

\subsection{Madeira lamelada colada}\label{ssec:material-mlc}

Adotou-se a classe GL24h, de composição homogênea, conforme a
EN 14080 \cite{EN14080:2013}. A \autoref{tab:gl24h} apresenta as propriedades
características empregadas no dimensionamento. A massa específica média é a
empregada no cálculo do peso próprio; a característica, nas verificações de
ligação.

\tabela{Propriedades características da classe GL24h}{\footnotesize
\begin{tabular}{lcrl}
\toprule
    \textbf{Propriedade} & \textbf{Símbolo} & \textbf{Valor} & \textbf{Unidade} \\
\midrule
    Resistência à flexão                      & $f_{m,k}$    & 24,0  & MPa \\
    Resistência à tração paralela             & $f_{t,0,k}$  & 19,2  & MPa \\
    Resistência à compressão paralela         & $f_{c,0,k}$  & 24,0  & MPa \\
    Resistência ao cisalhamento               & $f_{v,k}$    &  3,5  & MPa \\
    Resistência à tração perpendicular        & $f_{t,90,k}$ &  0,5  & MPa \\
    Resistência à compressão perpendicular    & $f_{c,90,k}$ &  2,5  & MPa \\
    Módulo de elasticidade médio              & $E_{0,mean}$ & 11.500 & MPa \\
    Módulo de elasticidade característico     & $E_{0{,}05}$   &  9.600 & MPa \\
    Módulo de cisalhamento médio              & $G_{mean}$   &    650 & MPa \\
    Massa específica característica           & $\rho_k$     &    385 & kg/m\textsuperscript{3} \\
    Massa específica média                    & $\rho_{mean}$&    420 & kg/m\textsuperscript{3} \\
\bottomrule
\end{tabular}}{Elaborado pelo autor (2026) com base na EN 14080 \cite{EN14080:2013}}{gl24h}{}{}

A resistência à tração perpendicular às fibras, de apenas 0,5 MPa, corresponde
a cerca de dois por cento da resistência à flexão. É essa desproporção que
torna determinantes as verificações da zona de cumeeira e do fendilhamento nas
ligações, discutidas no \autoref{cap:dimensionamento} e no
\autoref{cap:detalhamento}.

\subsection{Valores de cálculo}\label{ssec:valores-calculo}

O Eurocódigo 5 obtém o valor de cálculo de uma propriedade de resistência a
partir do valor característico por meio da expressão

\begin{align}
    X_d = k_{mod}\,\frac{X_k}{\gamma_M} \label{eq:valor-calculo}
\end{align}

\noindent
em que $\gamma_M$ é o coeficiente parcial do material, igual a 1,25 para
elementos de MLC e a 1,30 para ligações, e $k_{mod}$ é o coeficiente de
modificação.

O coeficiente $k_{mod}$ concentra duas influências que a madeira sofre e que os
materiais isotrópicos usuais não apresentam na mesma intensidade: o teor de
umidade de equilíbrio e a duração do carregamento. Quanto mais úmido o ambiente
e quanto mais prolongada a ação, menor a resistência mobilizável.

A umidade é tratada por classes de serviço. A classe 1 corresponde a ambientes
internos aquecidos, com umidade de equilíbrio da ordem de 12\%; a classe 2, a
ambientes cobertos porém não aquecidos, com umidade até cerca de 20\%; e a
classe 3, a condições mais severas que as anteriores, típicas de peças expostas
à intempérie ou em edificações abertas. A duração é tratada por cinco
categorias, que vão da ação permanente, superior a dez anos, à ação
instantânea. A \autoref{tab:kmod-kdef} apresenta os valores aplicáveis à MLC.
Quando a combinação reúne ações de durações distintas, adota-se o $k_{mod}$ da
ação de menor duração entre as presentes.

\tabela
{Coeficientes de modificação e de fluência para madeira lamelada colada}
{\footnotesize
\begin{tabular}{ccccccc}
\toprule
    \textbf{Classe de} & \multicolumn{5}{c}{\textbf{$k_{mod}$ por classe de duração da ação}} & \textbf{$k_{def}$} \\
\cmidrule(lr){2-6}
    \textbf{serviço} & Permanente & Longa & Média & Curta & Instantânea & \\
\midrule
    1 & 0,60 & 0,70 & 0,80 & 0,90 & 1,10 & 0,60 \\
    2 & 0,60 & 0,70 & 0,80 & 0,90 & 1,10 & 0,80 \\
    3 & 0,50 & 0,55 & 0,65 & 0,70 & 0,90 & 2,00 \\
\bottomrule
\end{tabular}}
{Adaptado da EN 1995-1-1 \cite{EN1995-1-1:2004}, tabelas 3.1 e 3.2}
{kmod-kdef}
{}
{}

O coeficiente $k_{def}$ quantifica a fluência e é empregado nas verificações de
estado-limite de serviço para converter a flecha instantânea na flecha final. A
progressão de 0,60 na classe 1 até 2,00 na classe 3 mostra que a deformação
diferida de uma peça exposta pode triplicar a parcela instantânea, o que faz da
classe de serviço uma decisão de projeto com consequência direta sobre a altura
da seção.

\subsection{Classe de serviço e durabilidade}\label{ssec:classe-servico}

Por se tratar de edificação aberta nas quatro faces, com as peças estruturais
sujeitas a variações de umidade e à exposição eventual à chuva conduzida
pelo vento, adotou-se a classe de serviço 3. A consequência direta é dupla:
na \autoref{eq:valor-calculo}, os coeficientes $k_{mod}$ caem para a faixa de
0,50 a 0,90, conforme a duração da ação, e o coeficiente de fluência assume
$k_{def} = 2{,}00$, o maior previsto na norma.

Essa classificação também determina as exigências de durabilidade. A
\autoref{tab:materiais-complementares} reúne as decisões correspondentes e os
demais materiais empregados. O adesivo Tipo I é o único admitido para peças em
classe de serviço 3, e o tratamento preservativo não dispensa o detalhamento de
proteção construtiva, tratado no \autoref{cap:detalhamento}.

\tabela
{Materiais complementares e exigências de durabilidade}
{\footnotesize
\begin{tabular}{lll}
\toprule
    \textbf{Item} & \textbf{Especificação} & \textbf{Referência} \\
\midrule
    Classe de uso da madeira           & UC 3.1     & EN 335 \cite{EN335:2013} \\
    Adesivo estrutural                 & Tipo I     & EN 301 \cite{EN301:2017} \\
    Tratamento preservativo            & Exigido    & EN 351-1 \cite{EN351-1:2007} \\
    Aço das barras de contraventamento & S235       & EN 1993-1-1 \cite{EN1993-1-1:2005} \\
    Aço das chapas de ligação          & S275       & EN 1993-1-1 \cite{EN1993-1-1:2005} \\
    Aço dos chumbadores                & Classe 8.8 & EN 1993-1-8 \cite{EN1993-1-8:2005} \\
    Concreto da fundação               & C25/30     & EN 1992-4 \cite{EN1992-4:2018} \\
\bottomrule
\end{tabular}}
{Elaborado pelo autor (2026)}
{materiais-complementares}
{}
{}

\section{VALORES CARACTERÍSTICOS DECLARADOS}\label{sec:valores-declarados}

Os dados de entrada dividem-se entre valores calculados pelo programa e valores
declarados como premissa. Os primeiros, entre eles o peso próprio da estrutura,
são obtidos por integração da geometria e apresentados no
\autoref{cap:modelagem}. Os segundos constam da
\autoref{tab:valores-declarados} e dependem de confirmação em ficha técnica de
fornecedor ou de definição em etapa posterior de projeto. A sobrecarga
corresponde à categoria H, de coberturas não acessíveis salvo para manutenção, e
a carga concentrada não é concomitante com a distribuída nem com o vento.

\tabela{Valores característicos declarados como dados de entrada}{\footnotesize
\begin{tabular}{p{4.6cm}crlp{3.4cm}}
\toprule
    \textbf{Grandeza} & \textbf{Símbolo} & \textbf{Valor} & \textbf{Unidade} & \textbf{Origem} \\
\midrule
    Telha metálica trapezoidal      & $g_{tel}$ & 0,10 & kN/m\textsuperscript{2} & Ficha técnica \\
    Instalações e serviços          & $g_{ins}$ & 0,10 & kN/m\textsuperscript{2} & Estimativa de projeto \\
    Sobrecarga de cobertura         & $q_k$ & 0,40 & kN/m\textsuperscript{2} & EN 1991-1-1 \cite{EN1991-1-1:2002} \\
    Carga concentrada de manutenção & $Q_k$ & 1,00 & kN & EN 1991-1-1 \cite{EN1991-1-1:2002} \\
    Velocidade básica do vento      & $V_0$ & 45,0 & m/s & ABNT NBR 6123 \cite{NBR6123:2023} \\
    Fator topográfico               & $S_1$ & 1,00 & -- & ABNT NBR 6123 \cite{NBR6123:2023} \\
    Fator estatístico               & $S_3$ & 1,00 & -- & ABNT NBR 6123 \cite{NBR6123:2023} \\
    Categoria de rugosidade         & -- & IV & -- & ABNT NBR 6123 \cite{NBR6123:2023} \\
    Classe da edificação            & -- & B & -- & ABNT NBR 6123 \cite{NBR6123:2023} \\
    Força de impacto de veículo, frontal & $F_{d,x}$ & 50,0 & kN & EN 1991-1-7 \cite{EN1991-1-7:2006} \\
    Força de impacto de veículo, lateral & $F_{d,y}$ & 25,0 & kN & EN 1991-1-7 \cite{EN1991-1-7:2006} \\
\bottomrule
\end{tabular}}{Elaborado pelo autor (2026)}{valores-declarados}{}{}

%% TODO: dado ausente — a velocidade básica $V_0 = 45$ m/s e o fator estatístico
%% $S_3 = 1{,}0$ constam como valores adotados no programa, mas a leitura da
%% isopleta correspondente ao local da obra e a justificativa do grupo de
%% classificação para $S_3$ não figuram na saída de cálculo. Ambos precisam ser
%% amarrados ao endereço definitivo da edificação.

%% TODO: dado ausente — as forças de impacto de 50 kN e 25 kN correspondem à
%% prática europeia para parques de veículos ligeiros e são, elas próprias,
%% dependentes de anexo nacional (ver checks/elu/impacto_veiculo.py). Na
%% ausência de anexo brasileiro, confirmar os valores com o orientador.

\section{PREMISSAS E RESERVAS NORMATIVAS}\label{sec:premissas-reservas}

\subsection{Ausência de anexo nacional}\label{ssec:ndp}

A adoção dos Eurocódigos fora do território europeu esbarra na inexistência de
anexo nacional brasileiro. Diversos parâmetros desses documentos são deixados à
determinação de cada país, e o texto do próprio Eurocódigo apenas recomenda um
valor. Na ausência de anexo, o projeto adota integralmente os valores
recomendados, e essa escolha é declarada de forma explícita em toda saída de
cálculo, conforme reproduzido na \autoref{tab:ndp}. São dezoito ao todo, e
caso a obra venha a ser submetida sob um anexo nacional concreto, todos devem
ser revistos.

\tabela
{Parâmetros de determinação nacional adotados}
{\scriptsize
\begin{tabular}{lrl}
\toprule
    \textbf{Parâmetro} & \textbf{Valor} & \textbf{Cláusula de origem} \\
\midrule
    $\gamma_G$ desfavorável            & 1,35  & EN 1990, tab. A1.2(B) \\
    $\gamma_G$ favorável               & 1,00  & EN 1990, tab. A1.2(B) \\
    $\gamma_G$ estabilizante, EQU      & 0,90  & EN 1990, tab. A1.2(A) \\
    $\gamma_G$ desestabilizante, EQU   & 1,10  & EN 1990, tab. A1.2(A) \\
    $\gamma_Q$                         & 1,50  & EN 1990, tab. A1.2(B) \\
    $\psi_0$ sobrecarga                & 0,00  & EN 1990, tab. A1.1, categoria H \\
    $\psi_1$ sobrecarga                & 0,00  & EN 1990, tab. A1.1, categoria H \\
    $\psi_2$ sobrecarga                & 0,00  & EN 1990, tab. A1.1, categoria H \\
    $\psi_0$ vento                     & 0,60  & EN 1990, tab. A1.1 \\
    $\psi_2$ vento                     & 0,00  & EN 1990, tab. A1.1 \\
    $\gamma_M$ madeira                 & 1,25  & EN 1995-1-1, tab. 2.3 \\
    $\gamma_M$ ligação                 & 1,30  & EN 1995-1-1, tab. 2.3 \\
    $k_{cr}$                           & 0,67  & EN 1995-1-1, item 6.1.7 \\
    Limite de flecha instantânea       & $L/300$ & EN 1995-1-1, tab. 7.2 \\
    Limite de flecha final líquida     & $L/250$ & EN 1995-1-1, tab. 7.2 \\
    $\alpha_{cc}$                      & 1,00  & EN 1992-1-1, item 3.1.6(1) \\
    $\gamma_C$                         & 1,50  & EN 1992-1-1, tab. 2.1N \\
    $\alpha_{cr}$ mínimo               & 10,0  & EN 1993-1-1, item 5.2.1(3), por analogia \\
\bottomrule
\end{tabular}}
{Elaborado pelo autor (2026)}
{ndp}
{}
{}

Coeficientes como $\gamma_M$ e os limites de
flecha incidem diretamente sobre o dimensionamento, e um anexo nacional que os
alterasse modificaria as taxas de utilização apresentadas no
\autoref{cap:dimensionamento}.

\subsection{Rota normativa para a ação do vento}\label{ssec:rota-vento}

A estrutura é uma cobertura isolada no sentido do item 7.2 da ABNT NBR 6123
\cite{NBR6123:2023}: uma cobertura sobre suportes de reduzidas dimensões, aberta
nas quatro faces. Os coeficientes da tabela 25, de coberturas isoladas a duas
águas, só se aplicam, porém, quando $0{,}07 \leq \operatorname{tg}\theta \leq 0{,}6$
e $h \geq 0{,}5\,\ell_2$ (item 7.2.1), em que $h$ é a altura livre entre o piso e a
aresta horizontal mais baixa da cobertura e $\ell_2$ a profundidade da cobertura
medida perpendicularmente à cumeeira. A primeira condição é atendida, com
$\operatorname{tg}\theta = \operatorname{tg} 10^{\circ} = 0{,}176$. A segunda não: adotou-se
$h = 7{,}00$ m, o pé-direito até o beiral, ao passo que $0{,}5\,\ell_2 = 20{,}25$ m
para a cobertura inteira, de 40,5 m, e ao menos 10,1 m para cada água dupla, de
20,25 m. A geometria está fora do domínio da tabela 25, e a norma não permite
extrapolá-la.

Para essa situação, e também quando obstruções possam ser colocadas sob a
cobertura, o último parágrafo do item 7.2.1 determina que a cobertura resista à
ação do vento, na zona de obstrução, calculada como edificação fechada de mesma
cobertura, com $c_{pi} = +0{,}8$ para obstruções na borda de sotavento e
$c_{pi} = -0{,}3$ para obstruções na borda de barlavento. As duas condições
ocorrem aqui, e os veículos podem ocupar toda a área coberta, de modo que a zona
de obstrução corresponde à cobertura inteira. Os coeficientes de pressão
externa são os da tabela 10 da norma, para telhados múltiplos, simétricos, de
tramos iguais, aplicável por ser $h = 7{,}00$ m inferior à largura do tramo,
$a' = 20{,}25$ m. Por permanecer a estrutura uma cobertura isolada, aplicam-se
ainda o item 7.2.2, à força de atrito com vento paralelo à cumeeira, e o item
7.2.6, aos elementos de vedação. O detalhamento e os padrões de coeficientes
varridos constam do \autoref{cap:modelagem}.

\subsection{Duração da ação do vento}\label{ssec:duracao-vento}

A EN 1995-1-1 \cite{EN1995-1-1:2004} admite classificar a ação do vento como de
curta duração ou como instantânea. Na classe de serviço 3, a primeira conduz a
$k_{mod} = 0{,}70$ e a segunda a $k_{mod} = 0{,}90$. Adotou-se a classificação como
curta duração, de resultado conservador, o que reduz em cerca de vinte e dois
por cento a resistência mobilizável nas combinações governadas pelo vento.
